\vsssub
\subsubsection{~$S_{is}$：依赖浮冰尺度的散射与耗散} \label{sec:IS2}
\vsssub

\opthead{IS2}{\ws}{F. Ardhuin, C. Sevigny, G. Boutin,  D. Dumont,  T. Williams}


该散射项的实现总体遵循 \cite{art:MM06} 的方法，并加入了对波浪破碎海冰的估计，以便更新最大浮冰直径。最后，还把一种基于蠕变的耗散同散射结合起来。


散射源项由散射系数 $\beta_{\mathrm{is},\mathrm{MIZ}}(\theta-\theta')$ 定义；默认情况下它是各向同性的，但可通过在 {\F SIS2} namelist 中设置 {\code IS2ISOSCAT=FALSE} 赋予任意方向依赖性。因此源项为
\begin{equation}
 \frac{S_{\mathrm{is}}(k,\theta)}{\sigma} =  \int_0^{2\pi}\beta_{\mathrm{is},\mathrm{MIZ}}(\theta-\theta') [s_\mathrm{scat}  N(k,\theta')-N(k,\theta)] d\theta'  .
 \end{equation}
其中 $s_\mathrm{scat}$ 默认设为 1.0，但可由 {\F SIS2} namelist 中的 {\code IS2BACKSCAT} 修改。\\


散射系数 $\beta_{\mathrm{is},\mathrm{MIZ}}$ 的确定基于理论反射系数 $\alpha_n(\sigma,h)$，该系数对应正入射波从开阔水域半平面传播到冰厚恒定为 $h$ 的覆冰水域半平面。由 \cite{art:KM08}（若 {\code IS2WIM1=0.}）或 \cite{art:Ben12}（若 {\code IS2WIM1=1.}）计算的 $\alpha_n(\sigma,h)$ 值制表存放在 {\code W3SIS2MD} 模块中。按照 \cite{art:Dea11}，破碎海冰被视为一系列这样的冰-水界面。忽略多次反射时，散射参数化把单位时间衰减定义为：网格单元中覆冰部分相当于由平均直径 $D_{m}$ 的一串浮冰组成，每块浮冰产生部分反射 $\alpha_n(\sigma,h)$，即
\begin{equation}
\beta_{\mathrm{is},\mathrm{MIZ}}=c_i~c_g \alpha_n(\sigma,h) / D_{m},
 \end{equation}
其中 $c_i$ 为冰密集度。\\
 

平均浮冰直径 $D_m$ 的估计基于如下假定：直径为 $D$ 的浮冰数量服从正比于 $D^{-\gamma}$ 的幂律。进一步假定该幂律适用于从最小值 $D_{\min}$ 到最大值 $D_{\max}$ 的 $D$ 范围。因此平均值为
\begin{equation}
D_m=\frac{\gamma}{\gamma -1}\times\frac{D_{\max}^{-\gamma +1}-D_{\min}^{-\gamma +1}}{D_{\max}^{-\gamma}-D_{\min}^{-\gamma}}.
\label{analytic_Dbar}
\end{equation}

\noindent
目前 $D_{\min}$ 为用户给定值。$D_{\max}$ 可作为强迫场提供，例如来自海冰模式或观测；或者当 namelist 参数 {\code IS2BREAK} 设为 {\code TRUE} 时，由局地波场对海冰的破裂估计得到。如果 namelist 参数 {\code IS2DUPATE} 设为 {\code TRUE}，即使波浪后来小到无法把海冰破裂到该尺度，较小的 $D_{\max}$ 仍会保持。在有外部强迫/耦合可用时（例如海冰模式中的海冰性质平流），这可能是合适的模式用法。相反，如果 {\code IS2DUPDATE} 设为 {\code FALSE}，$D_{\max}$ 将始终调整到局地海况，即使这意味着增大 $D_{\max}$。\\


假定波长为 $\lambda$ 的波浪破碎海冰后产生直径为 $\lambda / 2$ 的浮冰。在该参数化中，若满足以下三个条件，则发生波浪破冰 \citep{art:Wea13}：
\begin{enumerate}
\item $\lambda/2 \geq D_{\min} \quad \mathrm{and} \quad  \lambda/2 \leq D_{\max}$
\item $D_{\max}>D_c$，因为存在一个取决于海冰性质的临界直径，低于该直径时不可能发生弯曲破坏
\item $\varepsilon>\varepsilon_c$，入射波导致的应变必须大于给定临界应变
\end{enumerate}
第一个条件只是检查新的 $D_{\max}$ 值将大于 $D_{\min}$ 且小于先前的 $D_{\max}$。\\

第二个条件基于 \cite{inc:M86}，并采用 \cite{art:Bea18} 给出的修正，其中 $D_{c}$ 定义为
\begin{equation}
D_{c}=\dfrac{1}{2}\left(\frac{\pi^4 Y^* h^3}{48 \rho g (1 -\nu ^2)}\right)^{1/4}. 
\end{equation}

第三个条件对应弯曲应变阈值。波浪造成的水平应变与冰层曲率有关，在一维情况下为 $\varepsilon = 0.5 h \partial ^2 \eta_{\mathrm ice} / \partial x^2$。应变方差为
\begin{equation}
\langle\varepsilon^2\rangle = \left( \frac{h}{2} \right) ^2   \int_{k_1}^{k_2} k_{ice}^4 F(k)dk ,
\label{strain_sign}
\end{equation}
其中 $h$ 为冰厚，$k_{ice}$ 为波数 $2 \pi /\lambda_{ice}$。借鉴波浪破碎思想 \citep{art:BBY00}，曲率方差的积分限制在局地波数 $k_{\mathrm ice}$ 附近。还需注意，我们定义了一个有效最小冰厚 $h_{\min}$，使得若 $h<h_{\min}$，应变方差用 $h=h_{\min}$ 计算，以避免模式中出现与通常观测不符的“不可破碎”弹性薄冰。因此，我们把 $D_{\max}$ 取为满足以下判据的最短波的半波长：
\begin{equation}
F_{break}\, \sqrt{\varepsilon^2} > \frac{\sigma_c}{Y^*} ,
\label{strain_crit}
\end{equation}
其中 $\sigma_c$ 为海冰弯曲强度。$F_{break}$ 是表示随机波作用的因子，可通过 {\F SIS2} namelist 参数 {\code IS2BREAKF} 调整。理论上它应取决于海冰受波浪强迫的持续时间；基于 500 个 Rayleigh 分布波中的典型最大值，取 $F_{break}=3.6$。因此 $F_{break} \, E_s$ 为随机波的最大应变。\\




非弹性或非完全弹性耗散被包含在该例程而不是 ICn 中，因为它强烈依赖浮冰尺度。按照 \cite{art:Bea18}，假定浮冰变形并非完全弹性，且波浪诱导的循环应变和应力会把波能耗散为热。冰的循环变形可能需要远大于重力位能的弹性能，但这仅在海冰未破碎时成立。因此，若使用波面高程谱 $E(k)$，在海冰破碎或重新形成时可能会给 $E(k)$ 引入很大变化。相反，我们更倾向于使用能量谱 $ R C_g E(k)/C_{g,\mathrm{ice}}$，其中 $R$ 为 \cite{art:Wad73} 引入的系数，即弹性能与重力位能之比。
%For monochromatic waves of amplitude $a$, the energy by unit surface of waves in ice is thus 
%$R= \frac{1}{2} \rho g a^2 C_g$. It results that, without reflection nor dissipation, the wave variance in ice is such as :
%\begin{equation}
%E_{ice}=\frac{E\,C_{g}}{Cg_{ice}\,R}
%\label{eq:RWad}
%\end{equation}
%Note that if creep is activated (hence if {\code IS2BREAKE} $ > 0$), eq.~\ref{eq:RWad} 
%is also applied for the scattering curvature variance computation in $S_{IS2}$. 
对未破碎海冰，$R$ 为
\begin{equation}
R=1+C_R\frac{4Y^*h^3\pi^4}{3\rho g \lambda^4(1-\nu^2)},
\label{R}
\end{equation}
这里需注意 \cite{art:Wad73} 使用 $2h$ 作为冰厚；$C_R$ 默认通过 namelist 参数 {\code IS2BREAKE} 设为 1.0，但也可设为零，以改用真实高程谱。该因子 $R$ 也应用于波浪破冰计算。\\

默认使用非完全弹性耗散项，它通过 namelist 参数 {\code IS2ANDISB} 设为 {\code TRUE} 激活。非完全弹性耗散对应由波浪相关循环应力诱发的位错振荡运动中耗散为热的能量。当海冰承受正弦应力时，这种行为会在应力-应变图中表现为滞回，如 \cite{art:Cea98} 研究的图 4 所示。该文还提出了海冰非完全弹性行为模型，其应力-应变关系可计算滞回产生的椭圆面积。非完全弹性耗散被线性化为 $S_{\mathrm{ane}}=-\alpha_{\mathrm{ane}} E_{ice}$。\cite{art:Bea18} 从 \cite{art:Cea98} 的单色应力模型推导出 $\alpha_{\mathrm{ane}}$：
\begin{equation}
\alpha_{\mathrm{ane}}= \dfrac{A}{6}\left(k_i^{2}\dfrac{Y^*}{(1-\nu^{2})\rho gG}\right)^{2}h^{3}\frac{C_g}{GC_{g,i}} F_{\rm broken},
\label{eq:alpha_ane}
\end{equation}  
其中 $F_{\rm broken}$ 是从未破碎海冰到破碎海冰的启发式平滑过渡，使波长远大于浮冰尺度的波浪耗散逐渐趋于零，因为在这种情况下海冰不发生变形，也不会造成波能耗散：
\begin{equation}
F_{\rm broken}=\tanh \left( {\frac{D_{\max}-C_{\lambda}\lambda_{ice}} {D_{\max}C_{smooth} } }\right).
\label{smooth}
\end{equation}
非完全弹性或非弹性耗散在更新 $D_{\max}$ 后计算。
%dissipation is full for
%$\lambda_{ice}\leq D_{\max}$ and disappears for $C_{\lambda}\lambda_{ice}< D_{\max}$.  
该平滑过渡中的两个参数 $C_{\lambda}$ 和 $C_{smooth}$ 默认分别设为 $0.5$ 和 $0.4$，这是依据南大洋波浪资料调整得到的。它们可通过 namelist 参数 {\code IS2CREEPD} 和 {\code IS2CREEPC} 修改。最后，$A$ 为
\begin{equation}
A=\dfrac{4}{3}\sigma \alpha_{d}~\delta D^{d}~\dfrac{1}{\exp(\alpha_{d}s)+\exp(-\alpha_{d}s)}, 
\label{eq:A_ane}
\end{equation}  
其中各项在本节末尾表格中说明。耗散量级由位错柔度松弛参数 $\delta D^{d}$ 控制。\\


另一种非弹性耗散可通过把 {\code IS2ANDISB} 设为 {\code FALSE} 激活。此时使用如下耗散：

采用冰流动定律
\begin{equation}
%\left(\frac{d\varepsilon}{dt}\right)_{ij}=\frac{\tau^{n-1}}{B^n}\sigma_{i,j}'
\left(\frac{d\varepsilon}{dt}\right)_{ij}=\frac{\tau^2}{B^3}\sigma_{i,j}',
\end{equation}
$B$ 为流动定律常数，是海冰温度的函数。使用 \cite{art:Cea98} 由实验室实验估计的归一化参数 $A=10^{11}$，并取均匀冰温 270~K，可得 $B=10^7$~s$^{1/3}$。体积耗散率为
\begin{equation}
%\frac{de}{dt}= | \sigma_{xx}^{n+1}/(2B)^{n} |.
\frac{de}{dt}= | \sigma_{xx}^{4}/(2B)^3 |.
\end{equation}
非弹性耗散被线性化为 $S_{\mathrm{ine}}=-\alpha_{\mathrm{ine}} E_{ice}$。系数 $\alpha_{\mathrm{creep}}$ 由 \cite{art:Wad73} 的单色公式改写而来，等于
\begin{equation}
%\alpha_{\mathrm{ine}}=0.05 B h^5 \left(\frac{Y^*}{2B(1-\nu^2)}\right)^{(n+1)} I_n k^{n+1} \frac{C_g^2}{\rho g C_{g_{ice}}R^2} F \int_{k_1}^{k_2} k_{ice}^4 E(k)dk 
\alpha_{\mathrm{ine}}=0.05 B h^5 \left(\frac{Y^*}{2B(1-\nu^2)}\right)^{4} I_3 k^4 \frac{C_g^2}{\rho g C_{g_{ice}}R^2} F_{\rm broken} \int_{k_1}^{k_2} k_{ice}^4 E(k)dk, 
\label{eq:alpha_creep}
\end{equation}
其中 $I_3=\frac{1}{\pi} \int_0^\pi \sin^{4}\beta d\beta$，$F_{\rm broken}$ 已在上文定义。计算细节见 \cite{art:Bea18}。

最后，下表回顾 IS2 中使用的各个模式参数。有些参数作为常数定义在 {\code W3IS2MD} 模块中，另一些可通过 {\F SIS2} namelist 调整。

{\centering
\begin{tabular}{l  c c  c}
Parameters                   & Symbol        &  namelist parameter &  default values \\
\hline
Minimum floe size            & $D_{\min}$    & N. A. & 20~m \\
Initial floe size            & $D_{init}$    & N. A. & 1000~m \\
Ice fragility                & $\xi $        & N. A. & $0.9$ \\
Ice density                  & $\rho_{ice}$  & N. A. & 922.5~kg~m$^{-3}$ \\
Effective Young Modulus      & $Y^{*}$       & N. A. & 5.49~GPa \\
Poisson Coefficient          & $\nu$         & N. A. & 0.3 \\
Flexural strength            &  $\sigma_c$   & N. A. & 0.27~MPa \\
Flow law parameter           &$n$            & {\code IS2CREEPN} & 3\\
Flow law parameter           &$B$            & {\code IS2CREEPB} &  10$^7$~s$^{1/3}$\\
Elastic energy correction    & $C_R$         & {\code IS2BREAKE} & 1.0 \\
Relax. of disloc. compliance & $\delta D^{d}$& {\code IS2ANDISD} & $2  \times 10^{-9}$~(Pa$^{-1}$)\\
%Dislocation density          & $\Delta_d$    & N. A. & 1.8$\times10^{9}$ (m$^{-2}$)\\
%Restoring stress term        & $K$           & N. A. & 0.07 (Pa) \\
%Orientation factor           & $\Omega$      & N. A. & $\pi^{-1}$ \\
Burgers vector               & $b$           & N. A. & 4.52$\times10^{-10}$ (m) \\
Drag term                    & $B$           & N. A. & $B_0 \exp(Q_v/(k_bT_K)$ (Pa.s)\\
 -                           & $B_0$         & N. A. & 1.205$\times10^{-9}$ (Pa.s)\\
 Boltzmann constant          & $k_B$         & N. A. & 8.617$\times10^{-5}$ (eV.K$^{-1}$) \\
Activation energy            & $Q_v$         & {\code IS2ANDISE} & 0.55 (eV) \\
 Peak broadening term        & $\alpha_{d}$  & N. A. & 0.54 \\
 Reduced variable            & $s$           & N. A. & $\log(\tau\omega)$ \\
 Relax. time of disloc.      & $\tau$        & N. A. & $B/K$ (s) \\
 Temperature of sea ice      & $T_K$         & N. A. & $268.15$ (K)\\
\hline
\end{tabular}
}


